% ============================================================
%  Section 5: Simulation Design
% ============================================================

\section{Simulation Design}

\subsection{Overview}

To evaluate the three draft mechanisms under study, we build a multi-agent
simulation in which each NBA franchise is represented by an agent that receives
the current standings, its remaining schedule, the lottery odds structure for
the active mechanism, and the reward parameters, and then decides how much
effort to exert at regular intervals throughout the season.
We implement two classes of strategic agent.
The first is a \emph{rational agent} that maximizes expected utility analytically
at each decision point.
The second is an \emph{LLM agent} in which the team is represented by
\texttt{claude-haiku-4-5}; the model receives a structured natural-language
prompt describing the game state and reward structure and returns a continuous
effort level together with its reasoning.
Neither agent type encodes a hardcoded tanking rule: strategies emerge
endogenously from the incentive environment.
This lets us test whether the theoretical properties of each mechanism survive
contact with adaptive, strategically sophisticated agents---a more demanding
standard than the rule-based simulations used in prior work
\citep{kazachkov2020, highley2026}.

\subsection{Formal Model}

\paragraph{Teams and skills.}
There are $n = 30$ teams indexed $i = 1, \ldots, 30$.
Each team $i$ has a true skill level $\alpha_i$ drawn i.i.d.\ from a
log-normal distribution $\alpha_i \sim \text{LogNormal}(0,\, 0.3)$,
then sorted in ascending order so that team~1 is initially the weakest.
This calibration produces a realistic spread of talent: roughly $\pm 30\%$
around the league mean.

\paragraph{Season structure.}
Each team plays $G = 82$ games against opponents drawn from a balanced
round-robin schedule.
At season's end, the top $P = 16$ teams by win total qualify for the playoffs.
The remaining $L = 14$ teams enter the draft lottery under whichever mechanism
is active.

\paragraph{Win probability.}
We model game outcomes with a ratio model that incorporates an effort floor
$\delta = 0.3$.
When team $i$ plays at effort level $e_i \in [0,1]$ and team $j$ plays at
$e_j \in [0,1]$, their effective strengths are
\[
    \tilde{\alpha}_k = \alpha_k \bigl(\delta + (1 - \delta)\, e_k\bigr),
    \quad k \in \{i, j\},
\]
and the probability that team $i$ wins is
\[
    \Pr[\text{$i$ beats $j$}]
    = \frac{\tilde{\alpha}_i}{\tilde{\alpha}_i + \tilde{\alpha}_j}.
\]
The floor $\delta$ captures the fact that teams cannot perfectly throw a game:
even at $e = 0$ a team still plays at $30\%$ of its true strength.
At full effort ($e = 1$), $\tilde{\alpha}_i = \alpha_i$ and the model reduces
to a standard Bradley--Terry pairwise comparison.

\paragraph{Agent effort and decisions.}
Rather than a binary tank/try choice, agents select an effort level from the
discrete grid $e \in \{0,\; 0.25,\; 0.5,\; 0.75,\; 1.0\}$.
Values below $0.5$ are counted as a \emph{tanking decision} for reporting
purposes.
To reduce the number of API calls and to model realistic mid-season strategy
adjustments, decisions are made at \emph{checkpoints} spaced every 10 games
(games 0, 10, 20, \ldots, 70), giving each team 8 decisions per season.
The chosen effort level is held constant until the next checkpoint.

\paragraph{Payoffs.}
Each team's utility depends on whether it qualifies for the playoffs and, if
not, on which draft pick it receives.
Making the playoffs yields a value of $V = 50$, calibrated so that a playoff
berth is worth roughly half of the top draft pick.
For a non-playoff team assigned draft pick $k \in \{1, \ldots, 14\}$, the
draft payoff follows an empirically calibrated declining schedule:
\[
    D(k) \in \{100,\; 65,\; 45,\; 30,\; 22,\; 17,\; 13,\; 10,\;
               8,\; 7,\; 6,\; 5,\; 4,\; 3\},
\]
for $k = 1, \ldots, 14$ respectively, reflecting real NBA draft-pick trade
values.
Picks outside the lottery (positions $15$--$30$, awarded to playoff teams in
reverse order of their record) have a default value of $2$.

\paragraph{Lottery odds.}
For the current NBA lottery and the bilevel mechanism, lottery odds are
assigned by a weighted draw over the 14 non-playoff teams using the
2019-reform ticket weights
\[
    w = [140,\; 140,\; 140,\; 125,\; 105,\; 90,\; 75,\; 60,\;
         45,\; 30,\; 20,\; 15,\; 10,\; 5],
\]
where $w_1$ corresponds to the team with the worst record (or worst lock-point
record for bilevel).
The lottery determines picks 1--4 by sequential weighted sampling without
replacement; picks 5--14 are awarded in reverse-standings order.
Under COLA, draft order is determined entirely by accumulated ticket counts
rather than by a random draw; the team with the most tickets earns pick~1.

\paragraph{Multi-season skill dynamics.}
Because we simulate $S = 50$ seasons, team skills evolve over time.
After each season, the skill of team $i$ is updated as
\[
    \alpha_i^{t+1}
    = \operatorname{clip}\!\Bigl(
        \alpha_i^t
        - r\bigl(\alpha_i^t - \bar{\alpha}\bigr)
        + \phi\Bigl(1 - \frac{2(k_i^t - 1)}{n - 1}\Bigr),\;
        \alpha_{\min},\; \alpha_{\max}
      \Bigr),
\]
where $r = 0.10$ is a mean-reversion rate toward the league average
$\bar{\alpha} = 1.0$, $\phi = 0.06$ is the draft-pick boost magnitude,
$k_i^t$ is the draft pick received in season $t$, and the clip bounds are
$\alpha_{\min} = 0.25$ and $\alpha_{\max} = 2.50$.
Pick~1 delivers the maximum positive boost $(+\phi)$; the median pick delivers
no boost; pick~$n$ delivers the maximum penalty $(-\phi)$.
This dynamic creates the feedback loop between tanking incentives and
competitive balance that the reform mechanisms aim to address.

\subsection{LLM Agent Design}

Each of the 30 teams can be represented by an LLM agent instantiated with
\texttt{claude-haiku-4-5}.
At each checkpoint, the agent receives a structured prompt containing:
\begin{enumerate}
  \item its current win--loss record and rank among all 30 teams,
  \item the number of games completed and remaining in the season,
  \item the current standings for all 30 teams,
  \item the playoff cutoff and its probability of qualifying under honest play,
  \item a description of the active lottery mechanism and the expected draft
        pick value it would receive at each possible finishing position, and
  \item the full payoff schedule $V$ and $D(\cdot)$.
\end{enumerate}
A system-level prompt (cached per agent instance to reduce API cost)
explains the league rules, the scoring structure, and asks the agent to
return a JSON object \verb|{"effort": <float in [0,1]>, "reasoning": <string>}|.
No behavioral heuristic or tanking rule is encoded; the agent must reason from
the payoff structure alone.
Agents do not observe each other's decisions before making their own; choices
are made simultaneously at each checkpoint and resolved together.
For COLA, agents additionally receive their current ticket stockpile and a
description of the carryover and penalty rules, enabling multi-period reasoning.

\paragraph{Rational-agent surrogate.}
Because LLM agent runs are expensive, we also implement a \emph{rational agent}
that maximizes expected utility analytically.
At each checkpoint, the rational agent evaluates EU over the discrete effort
grid $\{0, 0.25, 0.5, 0.75, 1.0\}$ using a pre-computed Monte Carlo EV table
(5{,}000 lottery simulations) and the same projected-rank heuristic described
in Section~\ref{sec:metrics}.
Expected pick values under COLA are uniform across all lottery ranks---the
mechanism's key incentive-compatibility property---so rational agents under
COLA find no benefit to tanking within a season.
The rational agent provides a reproducible, deterministic baseline for the
systematic experiments in Section~5.5; the LLM agent is used for qualitative
validation of whether the same strategic patterns emerge in a less constrained
agent.

\subsection{Baseline: Honest Equilibrium}

Before running the strategic conditions, we establish a baseline by simulating
50 seasons under each mechanism with all 30 teams represented by a
\emph{honest agent} that always exerts full effort ($e = 1$) regardless of
incentives.
This baseline serves two purposes.
First, it provides a reference distribution for competitive balance: the
Kendall tau distance between final standings and the true skill ranking
$(\alpha_1 \leq \cdots \leq \alpha_{30})$ under honest play is the best
achievable outcome for the draft mechanism, since no tanking distorts
standings.
Second, it isolates the effect of stochastic game outcomes on competitive
balance from the effect of strategic behavior.
Under honest play, all three mechanisms produce identical tanking rates (zero
by construction) and near-identical competitive balance, confirming that any
observed divergence in the strategic conditions is attributable to the
incentive structure rather than to differences in how mechanisms handle
ties or ticket accumulation.

\subsection{Experimental Conditions}

We run five experimental conditions, summarized in Table~\ref{tab:conditions}.
Conditions A through C hold the mechanism fixed at the current NBA lottery and
vary the proportion of strategic agents, tracing the unraveling of the honest
equilibrium as rational behavior spreads.
Conditions D and E replace the current lottery with a reform mechanism while
keeping all agents strategic, testing whether the theoretical guarantees of
each reform survive adaptive behavior.

\begin{table}[h]
\centering
\renewcommand{\arraystretch}{1.3}
\begin{tabular}{cllp{5.5cm}}
\hline
\textbf{Cond.} & \textbf{Mechanism} & \textbf{Agent type} & \textbf{Purpose} \\
\hline
A & Current NBA (2019) & All honest
  & Baseline: competitive balance under zero tanking \\
B & Current NBA (2019) & Mixed: 1 rational, 29 honest
  & Does a single strategic team unravel the honest equilibrium? \\
C & Current NBA (2019) & All rational
  & Equilibrium tanking level under the current system \\
D & Bilevel ranking \citep{kazachkov2020} & All rational
  & Does the bilevel guarantee hold against strategic agents? \\
E & COLA \citep{highley2026} & All rational
  & Does COLA's incentive compatibility hold against strategic agents? \\
\hline
\end{tabular}
\caption{Experimental conditions. ``Rational'' agents maximize expected
utility; ``honest'' agents always exert full effort.}
\label{tab:conditions}
\end{table}

Each condition is run for $S = 50$ seasons.
In the mixed-population experiment (Condition~B), the single rational agent is
assigned to the team with the highest team ID (initially one of the stronger
teams), following the implementation in our codebase.
We additionally run a \emph{mixed-population sweep} that varies the number of
rational teams in $\{0, 1, 5, 15, 30\}$ under the current NBA lottery to
characterize how the aggregate tanking rate scales with the proportion of
strategic actors, and a \emph{playoff-value sensitivity sweep} that varies $V$
over $\{10, 30, 50, 80, 120\}$ to identify the threshold at which playoff
incentives dominate lottery incentives.

\subsection{Metrics}
\label{sec:metrics}

We evaluate each condition on three metrics.

\paragraph{Tanking rate.}
The fraction of all agent decisions in a season in which the chosen effort level
satisfies $e < 0.5$.
Because each of the 30 teams makes 8 decisions per season, the denominator is
$30 \times 8 = 240$ decisions per season-run.
A higher tanking rate indicates that the mechanism creates stronger incentives
to underperform.

\paragraph{Competitive balance (Kendall tau distance).}
Following \citet{kazachkov2020}, we measure competitive balance as the
normalized Kendall tau distance between the final standings ranking and the
true skill ranking over all $n = 30$ teams:
\[
    K(\sigma, \tau)
    = \frac{\bigl|\{(i,j) : \sigma(i) < \sigma(j)
                  \text{ but } \tau(i) > \tau(j)\}\bigr|}{\binom{n}{2}},
\]
where $\sigma$ is the ranking by end-of-season win total (rank~1 $=$ most wins)
and $\tau$ is the ranking by true skill $\alpha_i$ (rank~1 $=$ highest skill).
$K \in [0,1]$, with $K = 0$ indicating that standings perfectly reflect true
talent and $K = 1$ indicating a complete reversal.
Because skills evolve across seasons via the draft-pick update rule, we compute
$K$ each season using the skill vector $(\alpha_i^t)$ at the \emph{start} of
that season (before skills are updated), which is recorded in the
\texttt{skill\_history} table.

\paragraph{Draft efficiency.}
The average draft pick received by teams in each initial-skill quartile.
Teams in the bottom quartile (weakest 25\% by initial $\alpha_i$) should
receive the lowest pick numbers (best picks) if the mechanism is working as
intended.
We report the mean pick number for each quartile $Q \in \{1, 2, 3, 4\}$
averaged over all seasons and seeds; lower values for $Q = 1$ indicate better
redistributive efficiency.